% Evita que LaTeX estire el espacio vertical para llenar páginas
\raggedbottom

% Formato del título de capítulos y secciones
% \titleformat{\chapter}[block]{\titlerule[0.8pt]\normalfont\huge\bfseries}{\thechapter.}{.5em}{\Huge}[{\titlerule[0.8pt]}]
% \titlespacing*{\chapter}{0pt}{-19pt}{25pt}
% \titleformat{\section}[block]{\normalfont\Large\bfseries}{\thesection.}{.5em}{\Large}

\titleformat{\chapter}[block]{\normalfont\huge\bfseries}{\thechapter.}{.5em}{\Huge}[{\titlerule[0.8pt]}]
\titlespacing*{\chapter}{0pt}{-19pt}{25pt}
\titleformat{\section}[block]{\normalfont\Large\bfseries}{\thesection.}{.5em}{\Large}

% Formato del código fuente con lstlisting
\lstset{
  basicstyle=\ttfamily,
  breaklines=true,
}

% Márgenes
\geometry{
    a4paper,
    margin=2.75cm
}

% Limite de profundidad del índice
\setcounter{tocdepth}{2}

% Indentación de párrafos
\setlength{\parindent}{1cm}

\renewcommand{\lstlistingname}{Extracto de código}
\renewcommand*{\lstlistlistingname}{Índice de extractos de código}

%------------ Caption debajo de las tablas
\DeclareCaptionLabelSeparator*{spaced}{\\[1.5ex]}
\captionsetup[table]{textfont=it,format=plain,justification=justified,
  singlelinecheck=false,labelsep=spaced,skip=0pt}

% Placeholder para figuras pendientes. Sustituir por \includegraphics cuando exista la captura.
\newcommand{\placeholderfigura}[2][0.55\textwidth]{%
    \fcolorbox{black!40}{gray!15}{%
        \begin{minipage}[c][#2][c]{#1}
            \centering
            \large\textit{[figura pendiente]}
        \end{minipage}%
    }%
}